% additional.tex
\section{Additional results}\label{res_ar}
\subsection{Heterogeneous effects}
An important question is whether the detected anti-competitive pattern is uniform across all types of auctions or concentrated in particular contexts. We explore several dimensions of heterogeneity to see if certain conditions amplify or mitigate the incumbency and backlog discontinuities. First, we split the sample by auction format. The RDD analysis conducted separately for sealed-bid tenders (\emph{convite}) and for open descending auctions (\emph{pregão}) shows that the discontinuities are present in both formats. The point estimates for incumbency and backlog jumps are of similar order of magnitude in each subsample, suggesting that the advantage enjoyed by narrow winners is a pervasive phenomenon not limited to one auction mechanism. If anything, the incumbency effect is slightly larger in sealed-bid auctions, which could be consistent with the idea that less transparent, one-shot bidding environments might be more prone to favoritism. However, the difference is not statistically significant, and even in the more transparent descending auctions we still observe a significant incumbency and backlog gap at the threshold.

We also examine heterogeneity by contract size. We classify auctions as "small" or "large" based on the procurement's reference price relative to the median in our sample. The results indicate that both small and large procurements exhibit the incumbency advantage for narrow winners, but the effect on the backlog appears more pronounced in larger auctions. One interpretation is that higher-stakes contracts may attract greater effort by corrupt actors to favor certain firms, thereby producing a larger discrepancy in how much work the favored firm is handling relative to its competitors. On the other hand, smaller auctions, while still showing signs of favoritism, might not yield as large a backlog gap simply because the absolute values involved are lower. Additionally, we verified that the pattern holds across different categories of goods and services and across various procuring entities: in each major subgroup we considered, narrow winners consistently outperformed narrow losers in terms of past wins and backlog. This widespread presence of the effect reinforces the view that the underlying favoritism is not confined to a particular niche of procurement.

\subsection{The role of organized crime (PCC)}
Organized criminal networks can sometimes infiltrate public contracting, either by coercing legitimate bidders or by colluding with corrupt officials. In the context of São Paulo, the \emph{Primeiro Comando da Capital} (PCC) is a prominent criminal organization known to influence various illicit markets. We investigate whether areas or sectors with a higher suspected presence of the PCC exhibit different patterns of competition in procurement. To do this, we augment our analysis with proxy measures of organized crime influence, such as the incidence of violent crime or known PCC activity in the municipality of the purchasing unit. We then interact these proxies with our RDD framework to see if the incumbency or backlog discontinuities at the auction threshold are significantly altered in high-PCC-influence settings.

Our findings do not indicate a strong link between organized crime presence and the magnitude of the discontinuities. The incumbency and backlog jumps remain significant and of similar size even in regions where the PCC is believed to be active. In statistical terms, the interaction terms between the PCC proxy and the treatment effect at the threshold are small and not robustly significant. This suggests that the anti-competitive advantage enjoyed by narrow winners is not primarily driven by direct interference from organized crime cartels. Instead, the patterns we observe are more consistent with institutional or political factors (such as favoritism by public officials) rather than the influence of external criminal networks. In other words, while organized crime is a serious concern for public safety and possibly for certain high-value illegal markets, its role in routine public procurement favoritism appears limited according to our analysis.

\subsection{Correlation with predicted corruption}
If the favorable treatment of certain firms in procurement is a manifestation of corruption, one would expect it to correlate with other indicators of corrupt practices. We explore this by examining whether the degree of discontinuity we measure in an auction correlates with an independent measure of corruption risk. In practice, we construct a "predicted corruption" index at the level of each purchasing unit or municipality. This index could be based on past audit findings, the frequency of procurement irregularities, or other risk factors identified in the literature (for example, unusually high contract prices, low bidder turnout, or red-flag indicators in the tender process). 

When we relate our collusion indicators to this corruption index, we find a positive and significant correlation. Procurement environments that score higher on the predicted corruption scale tend to exhibit a larger incumbency advantage for narrow winners. In concrete terms, an auction conducted by an agency with a high corruption risk profile shows a more pronounced jump in incumbency at the threshold than an auction by a relatively clean agency. This relationship reinforces the interpretation that our RDD-detected pattern is indeed linked to underlying corrupt behavior. The favoritism signal we capture (narrow winners being consistently ahead in past wins) appears stronger in places or contexts that independent assessments would flag as vulnerable to corruption. This adds credibility to the argument that the observed discontinuities are not random or driven by benign factors, but align with scenarios of governance weakness where officials might be biasing outcomes. It also suggests that our method could serve as a complementary tool to traditional audit-based measures, helping to flag procurement processes that warrant further investigation.

\subsection{Firm outcomes}
Lastly, we consider the implications of these anti-competitive dynamics for the firms themselves and for the competitive landscape over time. If certain firms are systematically favored and keep winning repeatedly, one might expect those firms to grow larger or at least solidify their dominance in public procurement. Conversely, firms that consistently lose, especially in close competitions, might become discouraged or financially strained, potentially leading them to exit the market or reduce their participation in future tenders.

We examine the subsequent bidding behavior and success of firms, conditional on their experience in narrow auctions. The data suggest that firms which narrowly win an auction (and thereby benefit from the incumbency advantage) are more likely to win again in the near future and tend to bid in more auctions going forward. These favored firms often capture an increasing share of procurement spending over time, reinforcing their market position. On the other hand, firms that narrowly lose an auction appear less likely to win in their subsequent attempts compared to similar firms who had narrowly won. We observe a relative decline in participation among some of these consistently losing firms: they submit fewer bids in the months following a narrow loss and in some cases stop bidding altogether. This pattern is consistent with a discouragement effect, where firms that perceive the procurement process to be rigged in favor of certain competitors may reduce their effort or exit the public tender market.

The broader outcome of such dynamics is a less competitive supplier base and potentially higher prices for the government in the long run. As favored firms accumulate more contracts, their capacity and influence grow, which might further entrench their advantage in future bids. Meanwhile, the reduction in active competitors can lead to less pressure on prices and innovation. These findings underscore a potential vicious cycle: anti-competitive signals in procurement not only reflect past favoritism but can also shape the future market structure by squeezing out honest competition. This highlights the importance of timely detection and intervention to prevent long-term damage to the competitiveness of public procurement markets.
